% ============================================================================
% Introduction générale
% ============================================================================
\chapter*{Introduction générale}
\addcontentsline{toc}{chapter}{Introduction générale}
\markboth{Introduction générale}{}

Entre deux visites physiques, l'état d'un patient suivi pour une maladie chronique peut se dégrader sans que personne ne le voie~: une tension qui monte, une saturation en oxygène qui baisse, une fièvre qui s'installe. Le patient n'a pas toujours de moyen simple de signaler une mesure anormale, et le médecin n'a pas de visibilité continue sur ses patients. La télémédecine propose de réduire ce délai de détection en combinant deux usages complémentaires~: le \textbf{suivi à distance} de constantes physiologiques, qui signale automatiquement les valeurs inquiétantes, et la \textbf{téléconsultation} vidéo, qui permet de réagir immédiatement.

Ce mini-projet, deuxième sujet de la liste proposée par Prof.~Maria Zemzami~\cite{zemzami2026}, demande de \textbf{développer une application permettant la consultation vidéo et le partage de constantes physiologiques suivies dans le temps, avec un système d'alerte sur des seuils}. Trois livrables sont attendus~: un prototype fonctionnel (front et back), un scénario complet patient-médecin, et une analyse des enjeux réglementaires (consentement, confidentialité). Trois références normatives s'imposent~: la technologie WebRTC, le rapport de l'Organisation mondiale de la Santé (\ac{OMS}) sur la télémédecine, et la loi marocaine~131-13 qui encadre la télémédecine. Le cadre général du module ajoute la loi~09-08 sur la protection des données personnelles.

Nous avons réalisé \textbf{SahaLink}, un prototype complet développé à six et versionné dans un dépôt Git~: un serveur FastAPI avec une base SQLite, une interface web en français pour le patient et pour le médecin, un système d'alertes en temps réel par \ac{WS}, et une visioconsultation WebRTC en pair à pair. Toutes les données sont fictives. La sécurité n'est pas un chapitre décoratif~: elle est implémentée dans le code, testée automatiquement et reliée article par article aux textes de loi que nous avons lus.

Le rapport suit le plan du cahier des charges. Le chapitre~\ref{chap:etat} présente brièvement l'état de l'art. Le chapitre~\ref{chap:cadrage} cadre le contexte de soins, les acteurs, les besoins et les spécifications. Le chapitre~\ref{chap:conception} expose la conception (architecture, modèle de données, parcours, protocole), et le chapitre~\ref{chap:realisation} sa réalisation. Le chapitre~\ref{chap:validation} documente la validation et les résultats chiffrés, y compris le nombre de tests. Le chapitre~\ref{chap:securite} traite de la sécurité, de la confidentialité et de la conformité réglementaire, et le chapitre~\ref{chap:organisation} des aspects organisationnels. Le chapitre~\ref{chap:conclusion} dresse le bilan, les limites assumées et les perspectives.
